% Source PDF: source/普通高考/2010/2010北京理.pdf
% Page: 1; pdfimages -list was checked and no embedded image objects were available.
% Grid evidence: tmp/repaint_2010_beijing/source-1-grid.png (500px coarse, 100px fine)
% Original comparison crop: tmp/repaint_2010_beijing/q11_orig.png, cropped from the no-grid page render.
% Elements: frequency-density y-axis with arrow, broken x-axis, five adjacent height bins,
% horizontal dashed guides, class-boundary ticks, and frequency-density and 身高 labels.
% Relations: bins [100,110), [110,120), [120,130), [130,140), [140,150] have
% densities 0.005, 0.035, a (=0.030), 0.020, 0.010 respectively; each bin width is 10 cm.
% solid segments: axes, break mark, bar outlines, ticks and arrowheads.
% dashed segments: guides at 0.005, a, 0.020 and 0.010, ending at the corresponding bar boundary.
% labels: frequency-density above the y-axis, 身高 right of the x-axis arrow, tick labels below the axis.
\begin{tikzpicture}[baseline, line cap=round, line join=round]
\begin{axis}[
    width=12.0cm,
    height=5.4cm,
    axis lines=none,
    clip=false,
    xmin=-12,
    xmax=180,
    ymin=-0.004,
    ymax=0.041,
    ticks=none,
]
    % Frequency-density guides read from the original vector chart.
    \draw[dashed,line width=0.55pt] (axis cs:0,0.005) -- (axis cs:100,0.005);
    \draw[dashed,line width=0.55pt] (axis cs:0,0.030) -- (axis cs:120,0.030);
    \draw[dashed,line width=0.55pt] (axis cs:0,0.020) -- (axis cs:130,0.020);
    \draw[dashed,line width=0.55pt] (axis cs:0,0.010) -- (axis cs:140,0.010);

    % Five 10-cm height classes; a is determined by the stated total frequency.
    \addplot[
        ybar interval,
        draw=black,
        fill=white,
        line width=0.7pt,
    ] coordinates {
        (100,0.005)
        (110,0.035)
        (120,0.030)
        (130,0.020)
        (140,0.010)
        (150,0.000)
    };

    % Axes and the source's broken x-axis between 0 and 100 cm.
    \draw[->,line width=0.8pt] (axis cs:-1,0) -- (axis cs:160,0);
    \draw[->,line width=0.8pt] (axis cs:0,0) -- (axis cs:0,0.039);
    \fill[white] (axis cs:14,-0.0018) -- (axis cs:18,0.0018) -- (axis cs:22,-0.0018)
        -- (axis cs:26,0.0018) -- (axis cs:30,-0.0018) -- cycle;
    \draw[line width=0.8pt] (axis cs:14,-0.0018) -- (axis cs:18,0.0018)
        -- (axis cs:22,-0.0018) -- (axis cs:26,0.0018) -- (axis cs:30,-0.0018);

    \node[anchor=south west,inner sep=0pt,font=\small] at (axis cs:4,0.038) {\(\dfrac{\text{频率}}{\text{组距}}\)};
    \node[anchor=west,inner sep=0pt,font=\small] at (axis cs:166,0.001) {身高};

    \draw[line width=0.45pt] (axis cs:0,0) -- (axis cs:0,-0.0012);
    \node[anchor=north,inner sep=0pt,font=\scriptsize] at (axis cs:0,-0.0018) {0};
    \draw[line width=0.45pt] (axis cs:100,0) -- (axis cs:100,-0.0012);
    \node[anchor=north,inner sep=0pt,font=\scriptsize] at (axis cs:100,-0.0018) {100};
    \draw[line width=0.45pt] (axis cs:110,0) -- (axis cs:110,-0.0012);
    \node[anchor=north,inner sep=0pt,font=\scriptsize] at (axis cs:110,-0.0018) {110};
    \draw[line width=0.45pt] (axis cs:120,0) -- (axis cs:120,-0.0012);
    \node[anchor=north,inner sep=0pt,font=\scriptsize] at (axis cs:120,-0.0018) {120};
    \draw[line width=0.45pt] (axis cs:130,0) -- (axis cs:130,-0.0012);
    \node[anchor=north,inner sep=0pt,font=\scriptsize] at (axis cs:130,-0.0018) {130};
    \draw[line width=0.45pt] (axis cs:140,0) -- (axis cs:140,-0.0012);
    \node[anchor=north,inner sep=0pt,font=\scriptsize] at (axis cs:140,-0.0018) {140};
    \draw[line width=0.45pt] (axis cs:150,0) -- (axis cs:150,-0.0012);
    \node[anchor=north,inner sep=0pt,font=\scriptsize] at (axis cs:150,-0.0018) {150};
    \draw[line width=0.45pt] (axis cs:0,0.005) -- (axis cs:-2,0.005);
    \node[anchor=east,inner sep=0pt,font=\small] at (axis cs:-3,0.005) {0.005};
    \draw[line width=0.45pt] (axis cs:0,0.010) -- (axis cs:-2,0.010);
    \node[anchor=east,inner sep=0pt,font=\small] at (axis cs:-3,0.010) {0.010};
    \draw[line width=0.45pt] (axis cs:0,0.020) -- (axis cs:-2,0.020);
    \node[anchor=east,inner sep=0pt,font=\small] at (axis cs:-3,0.020) {0.020};
    \draw[line width=0.45pt] (axis cs:0,0.030) -- (axis cs:-2,0.030);
    \node[anchor=east,inner sep=0pt,font=\small] at (axis cs:-3,0.030) {$a$};
    \draw[line width=0.45pt] (axis cs:0,0.035) -- (axis cs:-2,0.035);
    \node[anchor=east,inner sep=0pt,font=\small] at (axis cs:-3,0.035) {0.035};
\end{axis}
\end{tikzpicture}
